\section{Conclusion}
\label{sec:conclusion}

\subsection{Reponse a la problematique}

La question posée était de savoir comment l'enchaînement d'une extraction répartie, d'une
zone de carrelage mouillé et d'un goulot unique altère le diagramme fondamental d'une foule
en panique, et lequel de deux mécanismes l'emporte : un effet tampon, où la traversée
désynchronise les arrivées et soulage la porte, ou un effet inertiel, où la faible adhérence
empêche les baigneurs de réguler leur approche. Dans notre modèle, la réponse tient en
trois constats.

\paragraph{Dans la piscine de référence, le carrelage mouillé n'altère rien de mesurable.}
Sur dix populations appariées, l'écart entre la piscine mouillée et son témoin sec vaut
$\Delta t_{50}=\SI{-0.01 \pm 0.07}{\second}$ (section~\ref{sec:e2}). Deux raisons
l'expliquent. La bande mouillée est à \SI{17}{\metre} de la porte : l'inertie s'efface sur le
sol sec avant que la file ne se forme. Surtout, la régulation vitesse--espacement garde les
baigneurs assez écartés pour qu'aucun n'ait à freiner fort. Mouiller la plage jusqu'à une
porte saturée ne suffit pas non plus, tant que l'adhérence reste celle d'un carrelage
mouillé ordinaire ($\Delta t_{100}=\SI{0.09 \pm 0.41}{\second}$).

\paragraph{Quand un sol très glissant touche la file d'une porte saturée, l'effet inertiel
l'emporte.} À \SI{1}{\metre\per\second\squared}, la porte de \SI{0.8}{\metre} perd
\SI{19}{\percent} de son débit et l'évacuation s'allonge de \SI{14}{\second} ; à
\SI{2}{\metre\per\second\squared}, de \SI{4}{\percent} et \SI{2.5}{\second}
(tableau~\ref{tab:e4-inertie}). Ces écarts sont positifs sur les dix graines de la campagne,
puis sur dix nouvelles graines fixées avant la relance, ce qui en fait le résultat le plus
solide de ce travail. Un témoin à la même adhérence, mais limitée à la bande de
\SI{3}{\metre}, ne change pas le débit : c'est bien le sol glissant sous la file qui agit. Le
diagramme fondamental mesuré devant la porte se déplace alors dans le sens de l'hypothèse
inertielle (figure~\ref{fig:diagramme-piscine}). La branche congestionnée monte jusqu'à
\num{3.9}~pers/\si{\metre\squared}, et à densité égale, entre \num{1.1} et
\num{1.7}~pers/\si{\metre\squared}, la file avance 20 à \SI{24}{\percent} moins vite. Notre
interprétation, qui reste à mesurer directement, est que ce n'est pas le freinage qui coûte,
puisque la projection arrête le baigneur contre celui qui le précède, mais la
\emph{relance} : à \SI{1}{\metre\per\second\squared}, retrouver \SI{1}{\metre\per\second}
depuis l'arrêt prend \SI{1}{\second} au lieu de \SI{0.34}{\second} sur le sec.

\paragraph{L'effet tampon n'a pas été mis en évidence.} Rien, dans nos campagnes de
référence, ne montre la traversée soulager la porte : quand le mouillé a un effet, il la
dégrade. L'extraction répartie, elle, compte : la porte ne limite l'évacuation que sous un
mètre de large, et au-delà, ce sont les franchissements et la population qui fixent la durée
(section~\ref{sec:e4}). Une étude exploratoire sans régulation laisse entrevoir des bouchons
plus fréquents sur sol sec que sur sol mouillé à faible panique, mais sur trop peu de graines
pour conclure : la question reste ouverte.

Notre réponse est donc conditionnelle : \emph{dans ce modèle, le carrelage mouillé n'altère
le diagramme fondamental que s'il est très glissant et qu'il touche la file d'une porte
saturée, et il l'altère alors par inertie, en rendant la file plus dense et plus lente}.

\subsection{Analyse critique}

\paragraph{Les adhérences qui produisent l'effet sont extrêmes.} \SI{1}{} et
\SI{2}{\metre\per\second\squared} (0,1 et 0,2~$g$) sont bien sous le seuil de glissance d'un
carrelage mouillé (0,34~$g$, Kendzior) : ce sont des sols savonneux. Le résultat dit qu'un
mécanisme inertiel existe dans le modèle et dans quelles conditions il apparaît, pas qu'un
carrelage de piscine ordinaire ralentit une évacuation.

\paragraph{Le modèle est étalonné sur un seul banc.} Le temps d'espacement
$T=\SI{0.35}{\second}$ ajuste le débit du goulot E0 à la mesure de Seyfried et al.\ sur sol
sec (\SI{2.8}{\percent} d'écart moyen, \num{5.7} et \SI{7.3}{\percent} sur les effectifs non
utilisés pour le choisir). Rien ne garantit qu'il convienne sur un sol glissant ou dans l'eau,
où nous l'appliquons faute de mieux. Plusieurs paramètres de l'eau et du mouillé sont
transposés de sources établies hors de leur cadre (une crue pour la vitesse dans l'eau, un
essai de glissance pour l'adhérence), et les durées de franchissement des échelles et du
rebord n'ont aucune source.

\paragraph{Le modèle est déterministe et sans contact mou.} Seul le placement initial est
tiré au hasard. Un bouchon qui forme une arche stable devant la porte ne cède donc jamais,
alors que les bouchons réels sont intermittents. L'exclusion dure par projection traduit un
choc parfaitement mou : il n'y a ni force de contact, ni poussée qui se transmettrait le long
d'une file. Les intensités que nous en reconstruisons dépendent du pas de temps et ne sont
que des ordres de grandeur. Le modèle est enfin bidimensionnel : nager et marcher dans l'eau
ne diffèrent que par la vitesse et l'adhérence.

\paragraph{Les trajectoires sur la plage sont peu réalistes.} Le guidage par le plus court
chemin fait marcher les baigneurs en file contre le bord du bassin. Une navigation locale
(Moussaïd et al.) corrige la marche, mais elle éloigne le banc E0 de la mesure de Seyfried ;
nous l'avons laissée hors de la référence plutôt que de refaire toutes les campagnes sans
l'avoir étalonnée.

\subsection{Pistes d'amelioration}

\begin{itemize}
  \item \textbf{Tester directement l'effet tampon} en étalant les départs par un temps de
        réaction à l'alarme, et en comparant le mouillé au sec à porte saturée.
  \item \textbf{Ajouter des fluctuations} au mouvement, pour que les bouchons puissent céder
        comme dans les expériences, et mesurer la durée des bouchons par la distribution des
        intervalles entre sorties plutôt que par $t_{100}$ seul.
  \item \textbf{Étalonner la navigation locale} sur le banc E0, comme la régulation, pour
        obtenir à la fois des trajectoires réalistes et le bon débit.
  \item \textbf{Mesurer les paramètres sans source} : chronométrer des sorties de bassin par
        l'échelle et par le rebord, et calibrer les taux de chute avant de les activer.
  \item \textbf{Mesurer les accélérations écrêtées} dans la file, pour confirmer que la
        relance, et non le freinage, porte l'effet inertiel.
\end{itemize}
